\documentclass[11pt]{article}
\usepackage{amsmath,amssymb}
\usepackage[margin=1in]{geometry}
\usepackage{hyperref}

\title{Per-species scaling of the spin-orbit coupling term in elkpy}
\author{}
\date{}

\begin{document}
\maketitle

\section{Setting: second-variational spin-orbit coupling in Elk}

Elk solves the Kohn-Sham problem in a second-variational scheme: a scalar-relativistic
Hamiltonian is diagonalized first, and spin-dependent terms (spin polarization,
spin-orbit coupling) are added afterward in the basis of scalar-relativistic
eigenstates. When \texttt{spinorb = .true.}, the spin-orbit term added inside each
muffin-tin sphere is the Koelling-Harmon (1977) approximation to the spin-orbit part of
the Dirac equation \cite{koelling1977}:
\begin{equation}
  \hat H_{\rm soc}(r) = f_{\rm soc}(r)\,\hat{\mathbf L}\cdot\boldsymbol\sigma,
  \qquad
  f_{\rm soc}(r) = \frac{1}{\bigl(2 M(r) c\bigr)^2}\,\frac{1}{r}\,
    \frac{\partial V_s(r)}{\partial r},
  \label{eq:hsoc}
\end{equation}
with the relativistic mass correction
\begin{equation}
  M(r) = 1 + \frac{1}{2c^2}\bigl(E - V_s(r)\bigr)\Big|_{E=0},
\end{equation}
where $V_s(r)$ is the spherical part of the Kohn-Sham effective potential in that
muffin-tin, $c$ is the speed of light (\texttt{solsc} in Elk's atomic-unit convention),
$\hat{\mathbf L}$ is the orbital angular momentum operator, and $\boldsymbol\sigma$ are
the Pauli matrices. This is exactly what Elk's \texttt{gensocfr.f90} computes per atom
(verified against \texttt{vendor/elk/src/gensocfr.f90}, whose own docstring cites the
same Koelling-Harmon reference).

Elk exposes one global scale on top of Eq.~\eqref{eq:hsoc}, the input parameter
\texttt{socscf} (default $1.0$): the radial coefficient actually evaluated is
\begin{equation}
  c_{\rm so} = \frac{y_{00}\cdot \texttt{socscf}}{4\,c^2},
  \qquad
  f_{\rm soc}(r) = \frac{c_{\rm so}}{r\,M(r)^2}\,\frac{\partial V_s(r)}{\partial r},
  \label{eq:socscf}
\end{equation}
the same scalar multiplying $\hat H_{\rm soc}$ for every atom in the cell, regardless
of species. Elk's manual states this knob's purpose plainly: it is used ``to enhance
the effect of spin-orbit coupling in order to accurately determine the magnetic
anisotropy energy (MAE)'' — i.e.\ it is a phenomenological correction factor, fitted
per material to compensate for whatever the scalar-relativistic/second-variational
approximation to Eq.~\eqref{eq:hsoc} misses relative to a full four-component Dirac
treatment or experiment, not a quantity derived from first principles.

\section{Why a single global scale is the wrong shape for multi-species cells}

$f_{\rm soc}(r)$ in Eq.~\eqref{eq:hsoc} is dominated by the region near the nucleus,
where $V_s(r)$ is steep; the resulting spin-orbit splitting grows strongly with atomic
number $Z$ (the same near-nuclear physics behind atomic fine-structure splitting,
heuristically $\sim Z^4$ for hydrogenic orbitals). A \texttt{socscf} value fitted to
correct the SOC-driven MAE of one species in a compound — typically a heavy element
whose scalar-relativistic treatment most understates true spin-orbit effects — is not,
in general, the correction a light species in the same cell needs. Because
\texttt{socscf} is a single number applied uniformly to every atom
(Eq.~\eqref{eq:socscf}), Elk upstream cannot express ``scale species A's spin-orbit
term but leave species B's alone.''

\section{elkpy's generalization}

elkpy promotes the scale in Eq.~\eqref{eq:socscf} from a single global number to a
per-species one, $\texttt{socscf} \to s_{\rm is}$, so that for an atom of species $is$:
\begin{equation}
  c_{\rm so}^{(is)} = \frac{y_{00}\cdot s_{is}}{4\,c^2}, \qquad
  s_{is} = \begin{cases}
    \texttt{soc\_scale[is]} & \text{if overridden}\\
    \texttt{socscf} & \text{otherwise (Elk's existing global default)}
  \end{cases}
  \label{eq:persp}
\end{equation}
Nothing about the physical content of Eq.~\eqref{eq:hsoc} changes — the radial
derivative $\partial V_s/\partial r$, the relativistic mass factor $M(r)$, and the
$\hat{\mathbf L}\cdot\boldsymbol\sigma$ structure are exactly Elk's existing
Koelling-Harmon term. Only the multiplicative prefactor becomes species-resolved, so a
fitted correction can be scoped to the species it was actually fitted for instead of
leaking onto every other species in the cell.

\subsection{Implementation}

\texttt{gensocfr.f90} already loops over every atom (\texttt{do ias=1,natmtot}) to
evaluate $f_{\rm soc}(r)$ for that atom's species, so
\texttt{patches/0001-per-species-soc-scale.patch} only inserts the species-dependent
lookup of Eq.~\eqref{eq:persp} inside that existing loop — it does not restructure the
term itself:
\begin{itemize}
  \item \texttt{modmain.f90}: a new array \texttt{socscfsp(maxspecies)}, sentinel
    $< 0$ meaning ``not overridden by elkpy, fall back to \texttt{socscf}''.
  \item \texttt{readinput.f90}: a new input block \texttt{elkpy\_socscale} (one
    \texttt{is~value} pair per line) that populates \texttt{socscfsp}.
  \item \texttt{gensocfr.f90}: inside the existing per-atom loop, \texttt{scso =
    socscf}, then \texttt{if (socscfsp(is) >= 0) scso = socscfsp(is)}, and $c_{\rm so}$
    is computed from \texttt{scso} instead of \texttt{socscf} directly.
\end{itemize}
This is the smallest possible hook into existing control flow, per the isolation
constraint in \texttt{docs/design.md} §8 — no new files were strictly needed here since
the per-atom loop already existed to hang the lookup on.

\subsection{Verification}

Verified against a real compiled Elk binary (\texttt{tests/test\_calculation\_soc.py}):
\begin{itemize}
  \item Single-species cell: \texttt{soc\_scale} reproduces the result of the
    equivalent global \texttt{socscf} value exactly (Eq.~\eqref{eq:persp} reduces to
    Eq.~\eqref{eq:socscf} when only one species is present).
  \item Two-species cell: scaling one species' \texttt{soc\_scale} changes that
    species' contribution independently of the other, confirming the override is
    applied per-species rather than falling back to a single effective global value.
\end{itemize}

\section{How to use in code}

\begin{verbatim}
from elkpy import Calculation

calc = Calculation(
    structure,
    spinpol=True,
    spinorb=True,
    soc_scale={"Fe": 1.5},   # scale only Fe's spin-orbit term; other species
                             # keep Elk's default global socscf (1.0 unless
                             # overridden separately via extra_blocks)
)
e = calc.get_energy()
\end{verbatim}
\texttt{soc\_scale} requires \texttt{spinorb=True} (elkpy raises otherwise, since a
scale on a disabled term has no physical effect), every key must name a species present
in the \texttt{Structure}, and every value must be $\geq 0$.

\begin{thebibliography}{9}
\bibitem{koelling1977}
D.~D.~Koelling and B.~N.~Harmon,
\emph{A technique for relativistic spin-polarised calculations},
J.\ Phys.\ C: Solid State Phys.\ \textbf{10}, 3107 (1977).
\end{thebibliography}

\end{document}
